\section{问题四的建模与求解}

本节围绕利润关键因素和未来预测目标，先筛选变量并构造时间特征，再比较候选模型，使用时间顺序验证选择最终模型，最后给出预测表和业务含义。

\subsection{问题分析}
说明预测目标、预测区间、可用信息和评价指标。

\subsection{变量筛选与特征工程}
说明变量角色、会计重构变量与运营变量的区别，以及缺失和极端值处理。

\subsection{候选预测模型与时间顺序验证}
\subsubsection{候选模型}
说明基线模型、线性模型和非线性模型的用途。

\subsubsection{滚动验证与评价指标}
说明训练窗口、预测步长、滚动起点和 MAE、RMSE、SMAPE 等指标。

\subsection{最终预测与误差带}
呈现模型比较、最终选择、逐日预测、合计预测和区间性质。

\subsection{业务启示}
说明预测对排班、备货、预算或资源配置的含义。

\subsection{问题四小结}
用简短结论收束，保留最终模型、预测结果和主要业务含义。
